% problem_id: S001-lemma-algebraization-principal-bis
% source_id: S001 (Stacks Project)
% source_file: algebraization.tex
% statement_locator: algebraization.tex lines 6518--6531
% proof_locator: algebraization.tex lines 6533--6654
% env: lemma
% id_note: label 在本来源内唯一
% pairing: tight (verified)
% status: complete (statement + author proof verbatim + reason + locators)
%
\begin{lemma}
\label{lemma-algebraization-principal-bis}
In Situation \ref{situation-algebraize} let $(\mathcal{F}_n)$ be an object
of $\textit{Coh}(U, I\mathcal{O}_U)$. Assume
\begin{enumerate}
\item $A$ is local and $\mathfrak a = \mathfrak m$ is the maximal ideal,
\item $A$ has a dualizing complex,
\item $I = (f)$ is a principal ideal,
\item $(\mathcal{F}_n)$ satisfies the $(2, 3)$-inequalities.
\end{enumerate}
Then $(\mathcal{F}_n)$ extends to $X$. In particular, if $A$ is
$I$-adically complete, then $(\mathcal{F}_n)$ is the completion
of a coherent $\mathcal{O}_U$-module. 
\end{lemma}

\begin{proof}
Recall that $\textit{Coh}(U, I\mathcal{O}_U)$ is an abelian category, see
Cohomology of Schemes, Lemma \ref{coherent-lemma-inverse-systems-abelian}.
Over affine opens of $U$ the object $(\mathcal{F}_n)$
corresponds to a finite module over a Noetherian ring
(Cohomology of Schemes, Lemma \ref{coherent-lemma-inverse-systems-affine}).
Thus the kernels of the maps $f^N : (\mathcal{F}_n) \to (\mathcal{F}_n)$
stabilize for $N$ large enough. By
Lemmas \ref{lemma-map-kernel-cokernel-on-closed} and
\ref{lemma-essential-image-completion}
in order to prove the lemma
we may replace $(\mathcal{F}_n)$ by the image of such a map.
Thus we may assume $f$ is injective on $(\mathcal{F}_n)$.
After this replacement the equivalent conditions of
Lemma \ref{lemma-equivalent-f-good} hold for the inverse system
$(\mathcal{F}_n)$ on $U$. We will use this without further mention
in the rest of the proof.

\medskip\noindent
We will check hypotheses (a), (b), and (c) of
Lemma \ref{lemma-when-done}. Hypothesis (b) holds by
Cohomology, Lemma \ref{cohomology-lemma-topology-I-adic-f}.

\medskip\noindent
Pick an inverse system of modules $\{M_n\}$ as in
Lemma \ref{lemma-system-of-modules}.
We may assume $H^0_\mathfrak m(M_n) = 0$ by replacing $M_n$ by
$M_n/H^0_\mathfrak m(M_n)$ if necessary. Then we obtain short exact
sequences
$$
0 \to M_n \to H^0(U, \mathcal{F}_n) \to H^1_\mathfrak m(M_n) \to 0
$$
for all $n$. Let $E$ be an injective hull of the residue field of $A$.
By Lemma \ref{lemma-helper-algebraize} and our current assumption (4)
we can choose, an integer $m \geq 0$, finite $A$-modules
$N_1$ and $N_2$ annihilated by $f^c$ for some $c \geq 0$ and
compatible systems of maps
$$
H^i_\mathfrak m(M_n) \to \Hom_A(N_i, E), \quad i = 1, 2
$$
for $n \geq m$
with the properties stated in the lemma.

\medskip\noindent
We know that $M = \lim H^0(U, \mathcal{F}_n)$ is an $A$-module whose
limit topology is the $f$-adic topology. Thus, given $n$, the module
$M/f^nM$ is a subquotient of $H^0(U, \mathcal{F}_N)$ for some $N \gg n$.
Looking at the information obtained above we see that
$f^cM/f^nM$ is a finite $A$-module. Since $f$ is a nonzerodivisor
on $M$ we conclude that $M/f^{n - c}M$ is a finite $A$-module.
In this way we see that hypothesis (c) of Lemma \ref{lemma-when-done} holds.

\medskip\noindent
Next, we study the module
$$
Ob = \lim H^1(U, \mathcal{F}_n) = \lim H^2_\mathfrak m(M_n)
$$
For $n \geq m$ let $K_n$ be the kernel of the map
$H^2_\mathfrak m(M_n) \to \Hom_A(N_2, E)$.
Set $K = \lim K_n$. We obtain an exact sequence
$$
0 \to K \to Ob \to \Hom_A(N_2, E)
$$
By the above the limit topology on $Ob = \lim H^2_\mathfrak m(M_n)$
is the $f$-adic topology. Since $N_2$ is annihilated by $f^c$
we conclude the same is true for the limit topology on $K = \lim K_n$.
Thus $K/fK$ is a subquotient of $K_n$ for $n \gg 1$.
However, since $\{K_n\}$ is pro-isomorphic to an inverse system of
finite length $A$-modules (by the conclusion of
Lemma \ref{lemma-helper-algebraize})
we conclude that $K/fK$ is a subquotient of a finite length
$A$-module. It follows that $K$ is a finite $A$-module, see
Algebra, Lemma \ref{algebra-lemma-finite-over-complete-ring}.
(In fact, we even see that $\dim(\text{Supp}(K)) = 1$ but
we will not need this.)

\medskip\noindent
Given $n \geq 1$ consider the boundary map
$$
\delta_n :
H^0(U, \mathcal{F}_n)
\longrightarrow
\lim_N H^1(U, f^n\mathcal{F}_N) \xrightarrow{f^{-n}} Ob
$$
(the second map is an isomorphism)
coming from the short exact sequences
$$
0 \to f^n\mathcal{F}_N \to \mathcal{F}_N \to \mathcal{F}_n \to 0
$$
For each $n$ set
$$
P_n = \Im(H^0(U, \mathcal{F}_{n + m}) \to H^0(U, \mathcal{F}_n))
$$
where $m$ is as above. Observe that $\{P_n\}$ is an inverse
system and that the map $f : \mathcal{F}_n \to \mathcal{F}_{n + 1}$
on global sections maps $P_n$ into $P_{n + 1}$.
If $p \in P_n$, then $\delta_n(p) \in K \subset Ob$
because $\delta_n(p)$ maps to zero in
$H^1(U, f^n\mathcal{F}_{n + m}) = H^2_\mathfrak m(M_m)$
and the composition of $\delta_n$ and $Ob \to \Hom_A(N_2, E)$
factors through $H^2_\mathfrak m(M_m)$ by our choice of $m$.
Hence
$$
\bigoplus\nolimits_{n \geq 0} \Im(P_n \to Ob)
$$
is a finite graded $A[T]$-module where $T$ acts via multiplication by $f$.
Namely, it is a graded submodule of $K[T]$ and $K$ is finite over $A$.
Arguing as in the proof of
Cohomology, Lemma \ref{cohomology-lemma-ML-general}\footnote{Choose
homogeneous generators of the form $\delta_{n_j}(p_j)$ for the displayed
module. Then if $k = \max(n_j)$ we find that for $n \geq k$
and any $p \in P_n$ we can find $a_j \in A$ such that
$p - \sum a_j f^{n - n_j} p_j$ is in the kernel of $\delta_n$
and hence in the image of $P_{n'}$ for all $n' \geq n$.
Thus $\Im(P_n \to P_{n - k}) = \Im(P_{n'} \to P_{n - k})$
for all $n' \geq n$.}
we find that the inverse system $\{P_n\}$ satisfies ML.
Since $\{P_n\}$ is pro-isomorphic to $\{H^0(U, \mathcal{F}_n)\}$
we conclude that $\{H^0(U, \mathcal{F}_n)\}$ has ML.
Thus hypothesis (a) of Lemma \ref{lemma-when-done}
holds and the proof is complete.
\end{proof}
